% !TEX root = chapter-standalone.tex

\section{Sum of Categories}
\label{sec:sum-of-categories}

\todotext{@J: review and edit}

The product of the previous section combines two categories by \emph{pairing} them.
The \emph{sum} of two categories combines them in a different, and in fact simpler, way: we place the two categories side by side, keeping both intact, and add nothing else.
An object of the sum is an object of one category \emph{or} of the other; morphisms within each part are as before; and there are \emph{no} morphisms connecting one part to the other.
Because the two parts remain disjoint, this construction is also called the \emph{disjoint union} of the two categories, and it is the categorical analogue of the disjoint union of sets.

\begin{ctdefinition}[Sum of categories]
    \label{def:disjoint-union-category}
    \SYNDEF{disjoint union of categories}
    Given two categories~\CatC and~\CatD, their \emph{sum} (or \emph{disjoint union})~$\CatC\Cplus\CatD$ is the category specified as follows:
    \begin{enumerate}
        \item \emph{Objects}: $\Ob_{\CatC\Cplus\CatD} \definedas \ObC \setdisunion \ObD$.
              That is, objects of $\CatC\Cplus\CatD$ are tuples of the form~$\disunionI{\Obja}$, with $i=1$ and $\Obja \setin \ObC$ or $i=2$ and~$\Obja \setin \ObD$.

        \item \emph{Morphisms}: Given objects~$\disunionI{\Obja}, \disunionJ{\Objb} \setin \Obof{\CatC\Cplus\CatD}$,
              \begin{equation}
                  \HomSet{\CatC\Cplus\CatD}{\disunionI{\Obja}}{\disunionJ{\Objb}} \definedas
                  \left\{
                  \begin{tabular}{ccc}
                      $\HomSet{\CatC}{\Obja}{\Objb}$ &  & if $i=j = 1$, \\
                      $\HomSet{\CatD}{\Obja}{\Objb}$ &  & if $i = j =2$, \\
                      $\Emptyset$                    &  & else.
                  \end{tabular}
                  \right.{}
              \end{equation}
              %
        \item \emph{Composition}: The composition operations $\mthenof{\CatC\Cplus\CatD}$, which are functions from
              $$
                  \HomSet{\CatC\Cplus\CatD}{\disunionI{\Obja}}{\disunionJ{\Objb}} \cartprod \HomSet{\CatC\Cplus\CatD}{\disunionJ{\Objb}}{\disunionK{\Objc}}
              $$
              to
              $$
                  \HomSet{\CatC\Cplus\CatD}{\disunionI{\Obja}}{\disunionK{\Objc}},
              $$
              are equal to
              $$
                  \left\{
                  \begin{tabular}{ccc}
                      $\mthenof\CatC$ &  & if $i=j = k= 1$, \\
                      $\mthenof\CatD$ &  & if $i = j = k =2$,
                  \end{tabular}
                  \right.{}
              $$
              and equal to the unique function
              $$\Emptyset \to \HomSet{\CatC\Cplus\CatD}{\disunionI{\Obja}}{\disunionK{\Objc}}$$
              in all other cases.
              %   that two morphisms $\mora$ and $\morb$ are composable only if they both  The composition~$\mora\colon \Obja \mto \Objb\setin \HomSet{\CatC + \CatD}{\disunionI{\Obja}}{\disunionI{\Objb}}$ and~$\morb\colon \Objb\mto \Objc\setin \HomSet{\CatC + \CatD}{\disunionJ{\Objb}}{\disunionJ{\Objc}}$, we have
              %       \begin{equation}
              %           \mora \mthenof{\CatC+\CatD} \morb \definedas
              %           \left\{
              %           \begin{tabular}{ccc}
              %               $\mora\mthenof\CatC \morb$  &  & if $i=j = 1$, \\
              %               $\mora \mthenof\CatD \morb$ &  & if $i = j =2$, \\
              %               does not exist              &  & else.
              %           \end{tabular}
              %           \right.{}
              %       \end{equation}
              %       \todotext{\alphubel: A little convoluted}
        \item \emph{Identity morphisms}: The identities  are copied from either category:
              \begin{align}
                  \catidofat{\CatC\Cplus\CatD}{\disunionA{\Obja}} & \definedas \catidofat\CatC\Obja, \\
                  \catidofat{\CatC\Cplus\CatD}{\disunionB{\Obja}} & \definedas \catidofat\CatD\Obja.
              \end{align}
    \end{enumerate}
\end{ctdefinition}

\begin{remark}
    If you think about categories in diagrammatic form, this operation corresponds to placing two categories side-by-side, without connecting them.
\end{remark}

\begin{remark}[Comparison with the sum of monoids]
    \label{rem:sum-cats-vs-sum-monoids}
    It is instructive to compare this construction with the sum of \SY{monoids} (\cref{def:sum-monoids}), which was considerably more involved: there, we had to interleave elements of the two \SY{monoids} into alternating words.
    The reason for the difference is that a \SY{monoid}, viewed as a category (\cref{def:monoidcat}), has exactly \emph{one} object, so a sum of \SY{monoids} must again have one object---the two structures are \emph{forced} to interact, and the interleaving bookkeeping is the price.
    Categories are under no such constraint: the sum of two categories can simply keep the two parts separate.
    In particular, note that
    \begin{equation}
        \monoidcat{\monoidA} \Cplus \monoidcat{\monoidB}
    \end{equation}
    is a category with \emph{two} objects, and is \emph{not} the categorification of the sum~$\monoidA + \monoidB$ (which has one object).
    Contrast this with the situation for products (\cref{exa:product-cat-monoids}), where the two constructions do agree.
\end{remark}

\begin{example}
    \label{exa:sum-cat-two-cities}
    Let~$\graphA$ and~$\graphB$ be the mobility graphs of two cities with no transport connection between them (say, the walking-and-subway networks of Zurich and Singapore), and consider the free categories of trips~$\pathcat{\graphA}$ and~$\pathcat{\graphB}$ (\cref{sec:catsfromgraphs}).
    The category~$\pathcat{\graphA} \Cplus \pathcat{\graphB}$ describes the combined system: a location is a location in one city or the other, trips within each city are as before, and there are no trips from one city to the other.
    An empty hom-set is the category theorist's way of saying ``you can't get there from here''.

    If we later add a connection (say, a flight), the combined model is no longer a plain sum---which is exactly the point: the sum is the special case of combination with \emph{zero} interaction.
\end{example}

\begin{example}
    \label{exa:sum-cat-posets}
    For \SY{posets}~$\posA$ and~$\posB$ viewed as categories (\cref{sec:posets-as-cats}), the sum~$\poscat{\posA} \Cplus \poscat{\posB}$ is the poset obtained by placing the two orders side by side, with elements of different parts incomparable.
    For instance, the sum of two chains $\bullet \mto \bullet \mto \bullet$ and $\bullet \mto \bullet$ is a five-element poset consisting of two independent chains.
\end{example}

\begin{remark}
    A remark similar to \cref{rem:n-fold-cart-products-of-cats} holds here, too: we can form the disjoint union of any finite number of categories.
    For example, $\CatA \Cplus \CatB \Cplus \CatC$ for three categories $\CatA, \CatB, \CatC$.
\end{remark}

\begin{exercise}
    Show that the \SYN{disjoint union of categories}{disjoint union} of two categories is indeed again a category.
\end{exercise}
\begin{solution}
    We check unitality and associativity.

    For unitality, consider a morphism~$\mora \setin \HomSet{\CatC\Cplus\CatD}{\disunionI{\Obja}}{\disunionI{\Objb}}$.
    (The index~$i$ is the same at source and target: by the definition of the hom-sets, no morphism connects objects of one part with objects of the other part.)
    If~$i = 1$, then~$\mora$ is a morphism of~\CatC, and
    \begin{equation}
        \catidofat{\CatC\Cplus\CatD}{\disunionI{\Obja}} \mthenof{\CatC\Cplus\CatD} \mora
        = \catidofat{\CatC}{\Obja} \mthenof{\CatC} \mora
        = \mora,
    \end{equation}
    using unitality in~\CatC; similarly~$\mora \mthenof{\CatC\Cplus\CatD} \catidofat{\CatC\Cplus\CatD}{\disunionI{\Objb}} = \mora$.
    The case~$i = 2$ is the same computation in~\CatD.

    For associativity, consider composable morphisms
    \begin{equation}
        \mora \setin \HomSet{\CatC\Cplus\CatD}{\disunionI{\Obja}}{\disunionI{\Objb}}, \quad
        \morb \setin \HomSet{\CatC\Cplus\CatD}{\disunionI{\Objb}}{\disunionI{\Objc}}, \quad
        \morc \setin \HomSet{\CatC\Cplus\CatD}{\disunionI{\Objc}}{\disunionI{\Objd}}.
    \end{equation}
    (Again, composability forces all the indices to agree.)
    If~$i = 1$, all three are morphisms of~\CatC, and both ways of composing them in~$\CatC\Cplus\CatD$ equal their composite in~\CatC, which is unambiguous by associativity in~\CatC:
    \begin{equation}
        (\mora \mthenof{\CatC} \morb) \mthenof{\CatC} \morc
        = \mora \mthenof{\CatC} (\morb \mthenof{\CatC} \morc).
    \end{equation}
    The case~$i = 2$ is analogous, with~\CatD in place of~\CatC.
\end{solution}
